\documentclass[12pt]{article}
\usepackage{graphicx} % Required for inserting images
\usepackage[margin=0.5in]{geometry}
\usepackage{amsmath, amssymb}
\usepackage{mathrsfs}


\usepackage{soul}


% Dr. David Brown Commands
\newcommand{\ds}{\displaystyle}
\newcommand{\R}{\mathbb{R}}
\newcommand{\N}{\mathbb{N}}
\newcommand{\Q}{\mathbb{Q}}
\newcommand{\C}{\mathbb{C}}


% Commands to get script r
\usepackage{calligra}
\DeclareMathAlphabet{\mathcalligra}{T1}{calligra}{m}{n}
\DeclareFontShape{T1}{calligra}{m}{n}{<->s*[2.2]callig15}{}
\newcommand{\scripty}[1]{\ensuremath{\mathcalligra{#1}}}

% Unit commands
\newcommand{\degree}{\text{°}}
\newcommand{\unitSpeed}{\frac{\text{m}}{\text{s}}}
\newcommand{\unitAcceleration}{\frac{\text{m}}{\text{s}^2}}

% Constant commands
\newcommand{\avogadro}{6.022\times10^{23}}

\newcommand{\boltzmannConstK}{1.381\times10^{-23}\frac{\text{J}}{\text{K}}}
\newcommand{\boltzmannConstInEV}{8.617\times10^{-5}\frac{\text{eV}}{\text{K}}}
\newcommand{\planckConst}{6.626\times10^{-34}\text{J}\cdot\text{s}}

\newcommand{\atmP}{1.01325\times10^5\,\text{Pa}}

% Commands to easily write partial derivatives
\newcommand{\partDeriv}[2]{\frac{\partial #1}{\partial #2}}
\newcommand{\doublePartDeriv}[2]{\frac{\partial^2 #1}{\partial^2 #2}}


\title{Problem 6.15}
\author{Gabriel Decker}


\begin{document}



\maketitle
\begin{flushleft}

I will compute the average energy of 10 atoms with the following energy distribution in two ways.
$0\,\text{eV}$ has 4 atoms.
$1\,\text{eV}$ has 3 atoms.
$4\,\text{eV}$ has 2 atoms.
$6\,\text{eV}$ has 1 atoms.\\~\\


\textbf{Part a:}\\
$$\overline{E}=\frac{4\cdot0\,\text{eV}+3\cdot1\,\text{eV}+2\cdot4\,\text{eV}+1\cdot6\,\text{eV}}{10}$$
$$\implies\overline{E}=\frac{3\,\text{eV}+8\,\text{eV}+6\,\text{eV}}{10}$$
$$\implies\overline{E}=\frac{17}{10}\,\text{eV}$$\\~\\




\textbf{Part b:}\\
The probabilities for each energy level is simply the multiplicity for each energy level. Therefore
\begin{table}[h]
    \centering
    \begin{tabular}{c|c}
        Energy Level (eV) & Probability\\
        \hline
        $0\,\text{eV}$ & $\frac{4}{10}$\\
        $1\,\text{eV}$ & $\frac{3}{10}$\\
        $4\,\text{eV}$ & $\frac{2}{10}$\\
        $6\,\text{eV}$ & $\frac{1}{10}$\\
    \end{tabular}
\end{table}\\~\\



\textbf{Part c:}\\
Using the following equation, we can calculate the average energy.
\begin{equation}
    \overline{E}=\sum_sE(s)\mathcal{P}(s)
\end{equation}

Plugging in our results from part b, we get the following.
$$\overline{E}=\left(0\cdot\frac{4}{10}+1\cdot\frac{3}{10}+4\cdot\frac{2}{10}+6\cdot\frac{1}{10}\right)\,\text{eV}$$
$$\implies\overline{E}=\left(\frac{3}{10}+\frac{8}{10}+\frac{6}{10}\right)\,\text{eV}$$
$$\implies\overline{E}=\frac{17}{10}\,\text{eV}$$\\~\\

Note that this is the same as in part a!


\end{flushleft}

\end{document}
